%% APPENDIX D -- the false-alarm expectation, the control ensemble, and the
%% execution-block quality criterion.  Owner: r11-chain.  Carries
%% \label{app:falsealarm} (04_method.tex, 08_backmatter.tex) and
%% \label{sec:blockquality} (04_method.tex).
\section{The false-alarm expectation and the control ensemble}
\label{app:falsealarm}

The paper's central claim is that the number of unattributed crossings is the
number chance predicts. This appendix gives that expectation, what it is
conditioned on, and what the control ensemble calibrating it can support.

\subsection{What the expectation is conditioned on}
%% (subsection label retired: nothing cites it)

The expectation is built from the control positions themselves, which are the
only empirical null the data provide, and it imposes the condition an event
must actually satisfy: to reach the trigger, and to fall outside the
attribution windows. Each of the \NWinA{} Class~A windows of the census
contributes the fraction of its own \NCtrl{} control positions that reaches
the trigger, which is the probability that a position in that field does so
--- no cell count enters it and no form is assumed for the tail. Summed over
the windows, \StExpTrig{} are expected to trigger by chance; the attribution
windows occupy \LgMaskOccPct{} per cent of the searched band; and the product
is \textbf{\StExpChance{} expected unattributed Class~A crossings against
\StObs{} observed}, the two sides taken over the same windows and under the
same condition.

The one assumption that cannot be read off the searched data is that a star
behaves like a control position in its own field, and the reserved blocks
measure that in the quantity used here. Their stars reach the trigger in
\StHoObs{} of \StHoNFine{} fine-channel windows where the control positions
of those same windows predict \StHoExpFine: a ratio of \StHoRatio, with
\StHoLo--\StHoHi{} allowed by resampling the reserved blocks, and no two of
the events in one block. In the \StHoNCoarse{} coarse-channel reserved
windows no star reaches the trigger, against \StHoExpCoarse{} predicted. No
departure is detected and none is applied, and the limit is the one so few
events allow: a departure smaller than about \StHoAccPct{} per cent would not
have been seen, and that is the accuracy of every chance expectation in this
paper.

Every Class~B crossing in the survey comes from the one execution block of
\S\ref{sec:blockquality}. Setting that block aside, the
coarse class adds almost nothing to either side --- the expectation over both
classes together is \StExpSurv{} against the same \StObs{} observed --- so
the Class~A comparison is also the survey's. Only \NCtrlWinObs{} of the
\NWindows{} windows hold any control position above
$T_{\rm trig}=\StageNullTrig$ at all, and \StageNullEligibleA{} of those are
Class~A: a coarse single-pointing window has too few cells for its own field
to reach the trigger by chance.

\emph{One set of windows was extracted twice, and the second extraction is
the one adopted.} \StNRep{} windows were first extracted over a field that
excluded integrations lying inside the primary beam; re-extracting them
recovers a median \StRepOnsrc{} times the on-source time and shifts the
search statistic by a median \StRepDT, or \StRepDTAbs{} in absolute value.
The rule for which integrations a window keeps was fixed before any statistic
in this paper was computed, and the re-extracted value is adopted wherever it
exists: it carries \StRepAdopted{} rows of Table~\ref{tab:ledger} and takes
\StRepFell{} crossings below the trigger. A re-extracted window's control
ensemble is re-extracted with it, so neither side of the comparison above
mixes the two. The deposited catalogue holds the
first extraction's statistic and the shorter species list in its own columns,
so its disposition column and the adopted ledger disagree on those
\StRepFell{} crossings and on \StRepAttr{} more that the species list now
attributes. The ledger is the authority for every disposition quoted here.

\subsection{Execution-block quality}
\label{sec:blockquality}

The comparison with the control positions assumes each block's ensemble
behaves like noise, and one block makes that assumption plainly false: every
one of its \NCtrl{} probes exceeds the trigger, the smallest at
\BqFailCtrlMin. A criterion for the condition was therefore fixed and applied
to every block, rather than recognised on the block that prompted it.

The statistic is the one this paper already uses as a diagnostic: the median
of a window's control ensemble, which is what separates a field with a source
in it from a field that is no null at all. That median is a maximum over the
window's own channel$\times$drift grid, so its level rises with the cell
count; it is therefore scaled by the median of the same quantity over the
windows of its own search class, and one threshold then means the same thing
in both. A block's statistic is the median of that ratio over its windows,
and \textbf{a block fails when it exceeds \BqFactor{} the scale of its
class} --- the factor already carried by the window-level data-quality flag,
fixed before that flag was applied to anything.

Applied to all \BqNBlocks{} blocks of the primary census, where the median
block sits at \BqQMed, the criterion fails \BqNFail: \BqFailBlock{} toward
\BqFailStar, at \BqFailQ{} where the next worst block is \BqNextQ. Nothing
else is close, and the unscaled form of the rule fails the same
\BqRawNFail{} block. Its \BqFailNWin{} windows are all coarse-channel and
hold the survey's entire Class~B crossing population, so the block carries
\EvNCrossFlag{} of the \EvNCrossRaw{} raw threshold crossings and no Class~A
statistic in this paper has ever contained it. \textbf{The population the
paper characterises is therefore the \EvNCrossQp{} crossings in
quality-passing data}: \EvNAttrQp{} attributed, \EvNUnattrQp{} unattributed,
\EvNRankQp{} of those outranking all \NCtrl{} control positions of their own
window against \EvNRankFlag{} in the failed block. The other
\EvNCrossFlag{} appear in Table~\ref{tab:ledger}, with their block named, and
enter no count anywhere else.

\subsection{The control ensemble: its calibration, and its resolution}
%% (subsection label retired: nothing cites it)

The ring geometry is set in \S\ref{sec:statistic}; its calibration and its
resolution are the subject here. Under that geometry the stellar ranks are
consistent with the uniform distribution exchangeability implies, in both
array strata and over the \VerNTested{} windows that retain a control vector,
the ACA windows measured on the geometry of each observation's own antenna
table and baselines rather than on a 12\,m beam. The test is not a powerful
one: a shift of less than about \VerKsPowerAca{} in the rank distribution
could not have been seen in samples of this size, so the strata are
consistent with uniform rather than demonstrably uniform.

\emph{The rank is not a significance test, and the reason is structural.}
Ranking the star against \NCtrl{} control positions drawn from the same field
is a good way to decide which windows deserve examination, but
$1/(\NCtrl+1)=\RankFloorVal$ is the \emph{resolution} of the ensemble --- the
smallest rank it can express --- and not the probability that noise produced
the window. A survey of \NWindows{} windows needs
$0.05/\NWindows=\BonfScaleCalc$, so that floor is \BonfRankRatio{} times too
coarse, and correlation widens it: neighbouring controls are correlated below
the synthesised-beam scale, so a window's annulus carries only
$N_{\rm eff}\simeq30$--$43$ independent spatial trials and the ensemble
resolves no better than \RankEffLo--\RankEffHi. Drawing more controls does
not rescue it, the primary beam bounding how many independent positions a
single pointing contains. \textbf{That floor is the design's central
limitation: no single window in this survey can be significant on the rank
alone, which is why confirmation here requires independent recurrence.}

\subsection{The trials budget}

The searched cells counted in \S\ref{sec:summary} are not independent
trials: channel smoothing and a drift grid stepped twice per channel
correlate neighbouring ones, which alone would leave \StGridFrac{} of them
independent. The control ensembles say otherwise. Their maxima sit at
\CtrlMaxRatioMed{} of the prediction for the maximum of the full searched
count treated as independent and Gaussian, a budget that predicts
\NCtrlWinExp{} windows carrying a control above $T_{\rm trig}$ against the
\NCtrlWinObs{} observed; and fitting each window's effective count against
its own ensemble returns \StNindFrac{} of its searched cells. Nothing quoted
here turns on which figure is right, because the expectation above is summed
from the measured exceedance rates and not from a cell count; by the
cell-count route it is \StExpModel, \StModelAgreePct{} per cent away.

\emph{The on-star crossings do exceed that budget, and the attribution
accounts for it.} A Gaussian tail over every cell the survey searched
predicts \CellExpect{} stellar exceedances against the \NCross{} found, but
\NAttributed{} of those are identified with a catalogued transition and
\EvNCrossFlag{} come from the block of \S\ref{sec:blockquality}; what is left
is the \StObs{} unattributed crossings in quality-passing data against
\StExpSurv{} predicted. No excess of chance crossings survives the
attribution, which is the comparison the paper rests on.

\emph{One stratum departs from uniformity and it is the one that matters.}
The quantity is the number of windows whose stellar peak both reaches the
trigger and exceeds all \NCtrl{} of that window's own control positions, and
over $n=\NFine$ fine-channel windows it is not uniform: \NStageOneFine{}
such windows against \ExpStageOneFine{} expected, while the coarse stratum,
three times larger and predicting \ExpStageOneCoarse, contains
\NStageOneCoarse{} --- including none of the \EvNCrossFlag{} crossings of the
flagged block, whose statistics stand below their own re-extracted control
maxima. A coarse channel dilutes a narrow circumstellar line, so
the class that can resolve one is the class that finds it --- and by the same
reasoning a genuine unresolved carrier would appear in Class~A first. The
departure survives removal of every fine window carrying a stellar crossing,
so fine-window rank probabilities are optimistic for the class as a whole and
not only where a crossing sits.

\emph{What is left is bounded, and the bounds are null.} Treating each
control as a pseudo-star and ranking it against the other \NCtrl-1 gives
\NPseudo{} ranks that the real stellar ranks agree with at
$\pv{\StarPseudoP}$; any systematic radial suppression of the controls is
below \RingRhoBoundPct{} per cent in the mean; and scrambling each window so
that nothing coherent survives along a drift track carries \LNNWindows{}
windows past the rank floor and changes no disposition.

%% ------------------------------------------------------------------
%% The chance expectation in this appendix is summed from the measured
%% control exceedance rates (stats_v413.py, round 260), so no cell count
%% and no Gaussian tail enters it; the cell-count route is quoted beside
%% it as a cross-check.  The rank-conditioned expectation
%% (\ClustStageMean against \ClustStageNObs) is a pipeline diagnostic and
%% must not be reachable as "the" expectation.
%% Floats held in the data release rather than here: tab:ctrlsep,
%% tab:excess, tab:arraysplit, fig:ctrldiag.  None is cited from the body.
%% ------------------------------------------------------------------

%% ★ tab:ledger MOVED HERE from Sec. 5 in the v4.12 length pass: the
%% crossing-by-crossing ledger belongs beside the false-alarm accounting that
%% explains how many of its rows are expected.  0.72 pp off the main text.
\begin{table*}
\centering
\caption{\textbf{The crossing ledger, from which every disposition in this
paper can be reproduced.} All \EvNCrossRaw{} threshold crossings, the
\EvNCrossFlag{} in the flagged block of \S\ref{sec:bdblock} among them and
labelled as such: the star, its execution block, the ALMA band, the crossing
frequency as the correlator delivered it and the same frequency in the star's
rest frame, the search statistic $T_\star$, the largest statistic
$T_{\rm ctrl}$ among the window's \NCtrl{} control positions, the nearest
masked transition and the stellar-frame offset $\Delta v_\star$ from it, then
the recurrence test: $n_{\rm rep}$ covering repeat blocks spanning $\Delta t$
days, the statistic $T_{\rm pers}$ a carrier persisting at the discovery flux
would have produced at the predicted stellar-frame cell, the statistic
$T_{\rm rep}$ measured there at the carried drift, combining the repeat
epochs by inverse variance, and the exclusion
$\sigma_{\rm excl}=(T_{\rm pers}-T_{\rm rep})/\sqrt{1+(T_{\rm pers}/T_\star)^2}$
(Eq.~\ref{eq:excl}), which is therefore reproducible from three columns of
this table. The fraction of the discovery amplitude still ruled out at the
trigger, $(T_{\rm rep}+5)/T_{\rm pers}$, follows from two of them.
$T_\star$ and $T_{\rm ctrl}$ always come from the same extraction of the same
window, so the comparison between them is one a reader can make;
$T_{\rm ctrl}$ gates nothing and says only which crossing was examined
first, and the drift-maximised reading of the cell is in
Table~\ref{tab:cand} for the two leading events and in the data release for
the rest. On the flagged rows $T_{\rm pers}$ and the exclusion, which would
inherit a discovery amplitude from the failing block, are withheld. A
crossing is attributed where
$|\Delta v_\star|\le\MaskHalfKms$\,km\,s$^{-1}$ and the transition named is
the nearest in that frame, so far from every transition it is the magnitude
and not the name that carries information. Block identifiers omit the
constant \texttt{A002\_} prefix.}
\label{tab:ledger}
\tiny
%% GENERATED by ledger_v410.py -- do not hand-edit.
%% One row per threshold crossing, one epoch convention, stellar-frame attribution.
%% part 1 of 2, 28 rows; the split is made in the generator so the float cannot outgrow its page.
\begin{tabular}{@{}l@{\hspace{0.9pt}}l@{\hspace{0.9pt}}c@{\hspace{0.9pt}}r@{\hspace{0.9pt}}r@{\hspace{0.9pt}}r@{\hspace{0.9pt}}r@{\hspace{0.9pt}}l@{\hspace{0.9pt}}r@{\hspace{0.9pt}}r@{\hspace{0.9pt}}r@{\hspace{0.9pt}}r@{\hspace{0.9pt}}r@{\hspace{0.9pt}}r@{\hspace{0.9pt}}l@{}}
\hline
Star & block & B & $\nu_{\rm topo}$ (GHz) & $\nu_{\star}$ (GHz) & $T_\star$ & $T_{\rm ctrl}$ & line & $\Delta v_{\star}$ & $n_{\rm rep}$ & $\Delta t$ (d) & $T_{\rm pers}$ & $T_{\rm rep}$ & $\sigma_{\rm excl}$ & disposition \\
\hline
BD+05 1668 & Xc7fa6f\_X28a8 & 6 & 240.1016 & 240.1054 & 73.398 & 74.69 & CS($5$--$4$) & -5911.9 & 29 & 0.9--113 & n/a & +0.15 & n/a & unattributed; flagged \\
BD+05 1668 & Xc7fa6f\_X28a8 & 6 & 224.7422 & 224.7458 & 48.303 & 83.74 & H$_2$CO($3_{1,2}$--$2_{1,1}$) & -1264.5 & 29 & 0.9--113 & n/a & +0.37 & n/a & unattributed; flagged \\
BD+05 1668 & Xc7fa6f\_X28a8 & 6 & 226.7422 & 226.7458 & 41.355 & 75.48 & CN($2_{3/2}$--$1_{1/2}$) & +88.0 & 29 & 0.9--113 & n/a & +0.66 & n/a & unattributed; flagged \\
BD+05 1668 & Xc7fa6f\_X28a8 & 6 & 242.1953 & 242.1992 & 39.586 & 83.29 & CS($5$--$4$) & -3349.2 & 29 & 0.9--113 & n/a & +0.77 & n/a & unattributed; flagged \\
$\beta$~Pic & Xf5d76d\_Xeb8 & 3 & 115.2603 & 115.2711 & 30.765 & 7.34 & CO($1$--$0$) & -0.2 & --- & --- & --- & --- & --- & attributed \\
$\beta$~Pic & Xf5d76d\_Xe19 & 3 & 115.2603 & 115.2711 & 27.684 & 6.48 & CO($1$--$0$) & -0.2 & --- & --- & --- & --- & --- & attributed \\
HD 48370 & Xc26103\_X155a & 6 & 230.5117 & 230.5243 & 26.834 & 18.36 & CO($2$--$1$) & -17.8 & --- & --- & n/a & --- & n/a & attributed; flagged \\
$\beta$~Pic & Xf4df6f\_Xc7 & 3 & 115.2620 & 115.2711 & 17.907 & 5.76 & CO($1$--$0$) & -0.2 & --- & --- & --- & --- & --- & attributed \\
$\beta$~Pic & Xf5d76d\_Xcc5 & 3 & 115.2604 & 115.2711 & 17.288 & 5.52 & CO($1$--$0$) & -0.2 & --- & --- & --- & --- & --- & attributed \\
$\beta$~Pic & Xf5d76d\_X32f1 & 3 & 115.2606 & 115.2713 & 14.654 & 7.27 & CO($1$--$0$) & +0.4 & --- & --- & --- & --- & --- & attributed \\
$\beta$~Pic & Xd9668b\_X3a90 & 6 & 230.5161 & 230.5378 & 11.681 & 9.12 & CO($2$--$1$) & -0.3 & --- & --- & --- & --- & --- & attributed \\
$\beta$~Pic & X6f1341\_X1484 & 6 & 230.5284 & 230.5379 & 9.832 & 8.20 & CO($2$--$1$) & -0.1 & --- & --- & --- & --- & --- & attributed \\
$\beta$~Pic & X7116f1\_X1785 & 6 & 230.5272 & 230.5384 & 9.677 & 14.13 & CO($2$--$1$) & +0.5 & --- & --- & --- & --- & --- & attributed \\
$\beta$~Pic & Xd9668b\_Xa9df & 6 & 230.5160 & 230.5377 & 8.948 & 8.41 & CO($2$--$1$) & -0.4 & --- & --- & --- & --- & --- & attributed \\
$\beta$~Pic & Xa7a216\_X23e4 & 6 & 230.5273 & 230.5377 & 7.988 & 6.03 & CO($2$--$1$) & -0.4 & --- & --- & --- & --- & --- & attributed \\
61 Vir & Xc079b5\_X82f & 7 & 346.3138 & 346.3232 & 5.958 & 5.65 & SO($9_{8}$--$8_{7}$) & -177.6 & 3 & 0.1--3.0 & 9.7 & +0.39 & 4.9 & unattributed; no recurrence \\
HD 285968 & Xd98580\_X3bf6 & 6 & 230.5021 & 230.5447 & 5.956 & 8.12 & CO($2$--$1$) & +8.7 & --- & --- & --- & --- & --- & attributed \\
HD 285968 & Xdb7ab7\_X5b4e & 6 & 230.5116 & 230.5447 & 5.868 & 7.22 & CO($2$--$1$) & +8.7 & --- & --- & --- & --- & --- & attributed \\
HD 285968 & Xdb4b9a\_X108f & 6 & 230.5100 & 230.5447 & 5.835 & 7.94 & CO($2$--$1$) & +8.7 & --- & --- & --- & --- & --- & attributed \\
CP$-$72 2713 & Xff0235\_X4a6d & 7 & 344.2697 & 344.2965 & 5.809 & 5.68 & SO($8_{8}$--$7_{7}$) & -12.3 & 5 & 0.1--2.1 & 6.3 & -0.82 & 4.8 & attributed \\
HD~139084~B & Xcd8029\_Xb6b0 & 7 & 345.8241 & 345.8241 & 5.763 & 17.70 & CO($3$--$2$) & +24.4 & --- & --- & --- & --- & --- & attributed \\
HD 14055 & Xfff3d8\_Xd46 & 7 & 331.7699 & 331.7664 & 5.575 & 5.99 & $^{13}$CO($3$--$2$) & +1068.7 & 30 & 0.9--650 & 170.4 & -0.02 & 5.6 & unattributed; no recurrence \\
HD 23484 & X10b22d7\_X2445 & 6 & 230.7904 & 230.8006 & 5.520 & 5.55 & CO($2$--$1$) & +341.5 & 4 & 0.1--2.1 & 10.4 & +1.25 & 4.3 & unattributed; no recurrence \\
HD 161868 & Xff45a6\_X10ca8 & 7 & 345.6338 & 345.6491 & 5.481 & 6.16 & SO($2_{3}$--$2_{1}$) & -48.0 & 13 & 0.0--6.1 & 92.2 & +0.24 & 5.5 & attributed \\
HIP 10679 & Xa95c04\_X14cd & 6 & 220.4185 & 220.4035 & 5.456 & 6.76 & $^{13}$CO($2$--$1$) & +6.5 & --- & --- & --- & --- & --- & attributed \\
HD~14055 & X11adad7\_X25c5e & 7 & 345.4564 & 345.4303 & 5.448 & 6.82 & SO($2_{3}$--$2_{1}$) & -237.8 & 30 & 0.1--653 & 14.0 & +0.03 & 5.1 & unattributed; no recurrence \\
HD 207129 & Xc19d6f\_X5706 & 6 & 230.4226 & 230.4070 & 5.433 & 5.52 & CO($2$--$1$) & -170.3 & 4 & 0.8--425 & 69.7 & -0.93 & 5.5 & unattributed; no recurrence \\
$\beta$~Pic & Xd9668b\_X3a90 & 6 & 241.7512 & 241.7739 & 5.420 & 5.81 & CS($5$--$4$) & -3869.7 & 1 & 1.0 & 5.7 & +0.72 & 3.4 & unattributed; no recurrence \\
\hline
\end{tabular}

\end{table*}

%% ★ GLENN, 2026-10-06: the bottom of this table ran into the text beneath
%% it.  A `table*` cannot break across pages and the ledger is a row per
%% threshold crossing, so the cure is structural and lives in the generator:
%% `ledger_v410.py` writes the body in parts of at most 28 rows and asserts
%% that it has enough floats for them.  This is the continuation; its caption
%% is one line, because the first carries the definitions.
\begin{table*}
\centering
\caption{\emph{The crossing ledger, continued.} The remaining
\CxLedgerPartTwo{} of the \EvNCrossRaw{} threshold crossings, in the same
order and with the same columns as Table~\ref{tab:ledger}, whose caption
defines them.}
\label{tab:ledgercont}
\tiny
%% GENERATED by ledger_v410.py -- do not hand-edit.
%% One row per threshold crossing, one epoch convention, stellar-frame attribution.
%% part 2 of 2, 28 rows; the split is made in the generator so the float cannot outgrow its page.
\begin{tabular}{@{}l@{\hspace{0.9pt}}l@{\hspace{0.9pt}}c@{\hspace{0.9pt}}r@{\hspace{0.9pt}}r@{\hspace{0.9pt}}r@{\hspace{0.9pt}}r@{\hspace{0.9pt}}l@{\hspace{0.9pt}}r@{\hspace{0.9pt}}r@{\hspace{0.9pt}}r@{\hspace{0.9pt}}r@{\hspace{0.9pt}}r@{\hspace{0.9pt}}r@{\hspace{0.9pt}}l@{}}
\hline
Star & block & B & $\nu_{\rm topo}$ (GHz) & $\nu_{\star}$ (GHz) & $T_\star$ & $T_{\rm ctrl}$ & line & $\Delta v_{\star}$ & $n_{\rm rep}$ & $\Delta t$ (d) & $T_{\rm pers}$ & $T_{\rm rep}$ & $\sigma_{\rm excl}$ & disposition \\
\hline
HR~1010 & Xc6fa08\_X1053 & 6 & 230.5133 & 230.5287 & 5.410 & 5.65 & CO($2$--$1$) & -12.0 & --- & --- & --- & --- & --- & attributed \\
ALMA J153702653-33192492 & Xff0235\_X8392 & 6 & 230.5222 & 230.5392 & 5.405 & 13.41 & CO($2$--$1$) & +1.6 & --- & --- & --- & --- & --- & attributed \\
HD 110411 & Xe45e29\_Xb7d4 & 6 & 219.4538 & 219.4357 & 5.393 & 5.85 & C$^{18}$O($2$--$1$) & -170.2 & 1 & 10 & 6.7 & +0.62 & 3.8 & unattributed; no recurrence \\
HD 31392 & X129754b\_X6946 & 6 & 231.1998 & 231.2231 & 5.289 & 5.81 & H($30\alpha$) & -876.3 & 8 & 0.1--41 & 15.7 & -0.26 & 5.1 & unattributed; no recurrence \\
HD 14055 & Xff99e1\_X24be0 & 7 & 331.7078 & 331.7035 & 5.277 & 5.69 & $^{13}$CO($3$--$2$) & +1011.6 & 30 & 0.9--652 & 252.2 & +0.15 & 5.3 & unattributed; no recurrence \\
AU Mic & Xa42f75\_X319 & 6 & 230.6864 & 230.6705 & 5.274 & 6.07 & CO($2$--$1$) & +172.3 & 2 & 310--455 & 6.4 & -0.16 & 4.2 & unattributed; no recurrence \\
HD 69830 & X12209f7\_X16edb & 7 & 322.4325 & 322.4444 & 5.259 & 5.42 & C$^{18}$O($3$--$2$) & -6268.5 & 5 & 0.1--29 & 8.9 & -0.48 & 4.8 & unattributed; no recurrence \\
ALMA J153702653-33192492 & Xfe62c1\_X398 & 6 & 216.1844 & 216.2036 & 5.228 & 6.02 & SiO($5$--$4$) & -1244.6 & 7 & 18--126 & 11.0 & +1.03 & 4.3 & unattributed; no recurrence \\
AU Mic & X899169\_X1838 & 6 & 230.8252 & 230.8294 & 5.221 & 5.85 & CO($2$--$1$) & +379.0 & 2 & 145--310 & 10.0 & +0.19 & 4.5 & unattributed; no recurrence \\
HD 14055 & Xff99e1\_X2c3bb & 7 & 345.5256 & 345.5216 & 5.186 & 5.66 & SO($2_{3}$--$2_{1}$) & -158.6 & 30 & 0.1--651 & 172.8 & +2.00 & 5.1 & unattributed; no recurrence \\
HD 14055 & Xfffde1\_X3835 & 7 & 345.0099 & 345.0070 & 5.177 & 6.56 & SO($2_{3}$--$2_{1}$) & -604.8 & 30 & 0.1--649 & 187.2 & -1.78 & 5.2 & unattributed; no recurrence \\
HD 14055 & Xff99e1\_X2b2ee & 7 & 344.7007 & 344.6965 & 5.175 & 5.92 & SO($8_{8}$--$7_{7}$) & +336.0 & 30 & 0.1--651 & 210.7 & -0.00 & 5.2 & unattributed; no recurrence \\
LHS 1140 & Xd9d7c7\_X6fd2 & 6 & 230.6355 & 230.6268 & 5.136 & 5.68 & CO($2$--$1$) & +115.5 & 4 & 5.0--27 & 7.3 & +0.33 & 4.0 & unattributed; no recurrence \\
HD 10647 & Xff294f\_Xaf70 & 7 & 331.6570 & 331.6943 & 5.120 & 5.91 & $^{13}$CO($3$--$2$) & +1003.3 & 10 & 0.1--592 & 125.8 & +0.39 & 5.1 & unattributed; no recurrence \\
CD$-$57 1054 & Xcf525f\_X5aa4 & 7 & 346.5099 & 346.5265 & 5.108 & 5.71 & SO($9_{8}$--$8_{7}$) & -1.7 & 1 & 39 & 6.0 & -0.27 & 4.1 & attributed \\
ALMA J153702653-33192492 & Xf8f6a9\_X10359 & 6 & 217.7292 & 217.7269 & 5.091 & 5.41 & H$_2$CO($3_{0,3}$--$2_{0,2}$) & -680.5 & 7 & 3.8--242 & 11.6 & +0.07 & 4.6 & unattributed; no recurrence \\
GJ 849 & Xd9668b\_X841a & 6 & 230.8769 & 230.8581 & 5.090 & 5.83 & CO($2$--$1$) & +416.3 & 5 & 0.1--51 & 9.3 & -1.69 & 5.3 & unattributed; no recurrence \\
ALMA J153702653-33192492 & Xfe62c1\_X398 & 6 & 218.8729 & 218.8923 & 5.090 & 5.56 & H$_2$CO($3_{2,1}$--$2_{2,0}$) & +181.2 & 7 & 18--126 & 10.5 & -1.04 & 5.0 & unattributed; no recurrence \\
HD 14055 & X1001fdc\_X29ae & 7 & 331.1559 & 331.1546 & 5.082 & 6.26 & $^{13}$CO($3$--$2$) & +513.8 & --- & --- & --- & --- & --- & unattributed; no spectrum \\
$\beta$~Pic & Xd9668b\_Xa9df & 6 & 241.8367 & 241.8594 & 5.079 & 5.99 & CS($5$--$4$) & -3765.1 & 1 & 1.0 & 4.8 & -0.63 & 3.9 & unattributed; no recurrence \\
ALMA J153702653-33192492 & Xf8f6a9\_X10359 & 6 & 218.6711 & 218.6688 & 5.048 & 5.77 & H$_2$CO($3_{2,1}$--$2_{2,0}$) & -125.1 & 7 & 3.8--242 & 11.6 & +0.00 & 4.6 & unattributed; no recurrence \\
HD 69830 & X12277ed\_X178c6 & 7 & 321.2005 & 321.2149 & 5.046 & 6.13 & C$^{18}$O($3$--$2$) & -7387.7 & 5 & 6.0--36 & 10.2 & -0.46 & 4.7 & unattributed; no recurrence \\
HR 1010 & Xc66ce1\_Xaa9 & 6 & 230.8348 & 230.8491 & 5.044 & 5.72 & CO($2$--$1$) & +404.6 & 7 & 1.1--14 & 10.3 & -1.53 & 5.2 & unattributed; no recurrence \\
61 Vir & Xc067f7\_X822 & 7 & 345.5599 & 345.5678 & 5.038 & 5.79 & SO($2_{3}$--$2_{1}$) & -118.5 & 3 & 2.0--3.0 & 9.3 & +0.45 & 4.2 & unattributed; no recurrence \\
HD 170773 & Xff0235\_X8c42 & 7 & 331.0101 & 331.0234 & 5.029 & 5.75 & $^{13}$CO($3$--$2$) & +394.8 & 4 & 0.9--113 & 74.8 & +1.56 & 4.9 & unattributed; no recurrence \\
ALMA J153702653-33192492 & Xd9d7c7\_Xbb64 & 7 & 347.2179 & 347.1890 & 5.009 & 5.68 & SiO($8$--$7$) & -122.2 & 1 & 0.1 & 5.0 & -2.06 & 5.0 & unattributed; no recurrence \\
HD~14055 & X1190de8\_Xd526 & 7 & 344.4676 & 344.4515 & 5.006 & 5.96 & SO($8_{8}$--$7_{7}$) & +122.7 & 30 & 20--611 & 11.6 & -0.60 & 4.8 & unattributed; no recurrence \\
$\eta$~Crv & X122b6ff\_X1041e & 7 & 357.2815 & 357.2475 & 5.006 & 5.50 & HCO$^{+}$($4$--$3$) & +431.3 & 1 & 87 & 11.0 & -0.62 & 4.8 & unattributed; no recurrence \\
\hline
\end{tabular}

\end{table*}
